\subsection{消融实验}
\label{sec:ablation}

为量化式~\eqref{eq:weight}中各因子的贡献，构造阶梯式消融序列：L0 为单像素基线（Coates 校正后匹配滤波取峰，无融合）；L1 在 L0 基础上对邻域做等权平均（无对齐）；L2 将等权替换为信号强度权重；L3 引入全轴匹配得分对齐但不用先验（搜索中心不依赖 $\Delta_{pq}$）；L4 先验整数对齐加几何软门控（仅 $g_\Delta$）；L5 加入匹配门与幅度响应（即式~\eqref{eq:weight}的完整权重但不含窗搜索）；L6 在 L5 基础上加入窗搜索重对齐（式~\eqref{eq:final_shift}，$K=5$）。各级均输出完整深度图并按同一口径评价，结果见表~\ref{tab:ablation230}与表~\ref{tab:ablation46}。

\begin{table}[!t]
\centering
\caption{230×310 分辨率下的消融实验结果（464 个场景全集；MAE 单位为 m，inl@1\% 为相对误差小于 1\% 的像素占比）}
\label{tab:ablation230}\scriptsize\setlength{\tabcolsep}{2.5pt}
\begin{tabular}{clcccccc}
\toprule
\multirow{2}{*}{级别} & \multirow{2}{*}{配置} & \multicolumn{2}{c}{$-5.9\pm8.7$} & \multicolumn{2}{c}{$-8.6\pm8.9$} & \multicolumn{2}{c}{$-15.4\pm9.4$} \\
\cmidrule(lr){3-4}\cmidrule(lr){5-6}\cmidrule(lr){7-8}
 & & MAE & inl@1\% & MAE & inl@1\% & MAE & inl@1\% \\
\midrule
L0 & 单像素（Coates+MF） & 0.716 & 39.8 & 0.860 & 33.5 & 1.727 & 13.2 \\
L1 & 不对齐 + 等权平均 & 0.084 & 62.6 & 0.099 & 62.0 & 0.199 & 57.7 \\
L2 & 不对齐 + 信号强度权重 & 0.082 & 62.7 & 0.098 & 62.1 & 0.196 & 58.0 \\
L3 & 无先验 + 全轴匹配对齐 & 1.571 & 29.5 & 1.653 & 25.4 & 2.055 & 12.2 \\
L4 & 先验整数对齐 + 几何门 & 0.064 & 64.3 & 0.080 & 63.3 & 0.194 & 57.1 \\
L5 & L4 + 匹配门 + 信号强度 & 0.061 & 65.1 & 0.077 & 64.0 & 0.190 & 57.7 \\
L6 & L5 + 窗搜索重对齐 & \textbf{0.058} & \textbf{64.3} & \textbf{0.074} & \textbf{63.9} & \textbf{0.186} & \textbf{57.6} \\
\bottomrule
\end{tabular}
\end{table}

\begin{table}[!t]
\centering
\caption{46×62 分辨率下的消融实验结果（464 个场景全集；MAE 单位为 m，inl@1\% 为相对误差小于 1\% 的像素占比）}
\label{tab:ablation46}\scriptsize\setlength{\tabcolsep}{2.5pt}
\begin{tabular}{clcccccc}
\toprule
\multirow{2}{*}{级别} & \multirow{2}{*}{配置} & \multicolumn{2}{c}{$-5.7\pm8.4$} & \multicolumn{2}{c}{$-8.4\pm8.6$} & \multicolumn{2}{c}{$-15.2\pm9.0$} \\
\cmidrule(lr){3-4}\cmidrule(lr){5-6}\cmidrule(lr){7-8}
 & & MAE & inl@1\% & MAE & inl@1\% & MAE & inl@1\% \\
\midrule
L0 & 单像素（Coates+MF） & 0.665 & 37.0 & 0.814 & 30.9 & 1.695 & 12.5 \\
L1 & 不对齐 + 等权平均 & 0.323 & 22.7 & 0.349 & 22.3 & 0.522 & 19.7 \\
L2 & 不对齐 + 信号强度权重 & 0.317 & 23.0 & 0.342 & 22.5 & 0.519 & 20.0 \\
L3 & 无先验 + 全轴匹配对齐 & 1.470 & 27.3 & 1.533 & 23.2 & 2.031 & 10.0 \\
L4 & 先验整数对齐 + 几何门 & 0.214 & 29.3 & 0.235 & 30.1 & 0.460 & 24.1 \\
L5 & L4 + 匹配门 + 信号强度 & 0.197 & 32.4 & 0.221 & 32.5 & 0.451 & 25.2 \\
L6 & L5 + 窗搜索重对齐 & \textbf{0.184} & \textbf{40.1} & \textbf{0.207} & \textbf{38.2} & \textbf{0.427} & \textbf{27.5} \\
\bottomrule
\end{tabular}
\end{table}

结果呈现三个一致的规律。其一，无先验的数据驱动对齐（L3）在两个分辨率下均系统性崩溃（MAE $1.5\sim2.1\,\mathrm{m}$，甚至劣于不融合的 L0）：归一化匹配得分在背景主导的直方图上频繁锁定背景巧合峰，且跨邻居的错位独立发生、无法通过平均抵消。将搜索中心约束为先验整数对齐（L4）后误差立即恢复一个数量级以上，验证了 \ref{sec:alignment} 节“单目先验提供搜索范围、数据决定最终对齐”的分工必要性；对照式~\eqref{eq:delta}，此时起作用的仅是单目相对位移的几何趋势，其全局尺度误差不参与。其二，几何门（L4）与匹配门、信号强度（L5）的逐级引入在两个分辨率下均带来单调改善，说明几何门、匹配门与幅度响应的贡献可分离且方向一致。其三，窗搜索（L6）在统一 $K=5$ 配置下于两个分辨率、三个照度档全部取得最优 MAE 与中位数误差，相对 L5 的净增益为 $3\%\sim7\%$；同时 inl@1\% 不降，表明小窗重对齐在修正邻居残余错位的同时未引入背景巧合峰。

值得指出的是 MAE 与 inl@1\% 在 L0$\to$L1 处的背离：46×62 上等权平均使 MAE 从 $0.665$ 降至 $0.323$，inl@1\% 却从 $37.0\%$ 跌至 $22.7\%$。这是误差分布重组的结果——单像素估计呈双峰分布（约 37\% 的像素亚厘米级准确，其余误差约 $1\,\mathrm{m}$ 量级），等权平均将高信噪比像素与跨边缘、饥饿邻居无差别混合，以头部精度换取尾部均值；后续几何门与匹配门逐级将该头部精度找回，至 L6 时 inl@1\% 反超 L0（$40.1\%$ vs $37.0\%$），实现两端兼顾。这一现象也说明单靠 MAE 会高估朴素融合的价值。
